\section{Data and Evaluation Protocol}
\label{sec:protocol}

\Cref{fig:workflow} summarises the study.
All models, in every setting, share one split, one model-selection rule and one test set, so that differences between rows of a table can be attributed to the model or to the training distribution and not to the evaluation.

\begin{figure}[t]
  \centering
  \resizebox{\columnwidth}{!}{% Experimental workflow (Fig. 1). Fill colours follow the reference categorical palette.
\definecolor{wfblue}{HTML}{2A78D6}
\definecolor{wforange}{HTML}{EB6834}
\definecolor{wfaqua}{HTML}{1BAF7A}
\definecolor{wfink}{HTML}{52514E}
\begin{tikzpicture}[
    font=\scriptsize, >=Stealth, node distance=1.6mm,
    box/.style={draw=wfink, rounded corners=1.5pt, line width=0.4pt, align=center, inner sep=1.5pt,
                minimum height=7mm, text width=#1},
    box/.default=18mm,
    grp/.style={draw=wfink!60, dashed, rounded corners=2pt, inner sep=2mm},
    gl/.style={font=\scriptsize\bfseries, text=wfink, anchor=base, yshift=1.2mm},
    arr/.style={->, wfink, line width=0.5pt}
  ]
  % 1. Splits
  \node[box=16mm] (pool) {Train pool\\544 / class};
  \node[box=16mm, below=of pool] (val) {Validation\\96 / class};
  \node[box=16mm, below=of val] (test) {Test\\160 / class};
  \node[grp, fit=(pool)(val)(test)] (G1) {};
  \node[gl] at (G1.north) {SkyView split};

  % 2. Training sets
  \node[box=18mm, right=7mm of pool, fill=wfblue!12] (bal) {Balanced\\544 / class};
  \node[box=18mm, below=of bal, fill=wforange!12] (lt) {Long-tailed\\544 $\to$ 34 / class};
  \node[box=18mm, below=of lt, fill=wfaqua!12] (rb) {Rebalanced\\LT + augmentation};
  \node[grp, fit=(bal)(lt)(rb)(bal |- G1.north)(rb |- G1.south), inner sep=0pt] (G2) {};
  \node[gl] at (G2.north) {Training sets};

  % 3. Models
  \node[box=21mm, right=7mm of bal] (hc) {Hand-crafted\\LBP / colour / SIFT};
  \node[box=21mm, below=of hc] (cnn) {CNNs\\ResNet-18, EffNet-B3};
  \node[box=21mm, below=of cnn] (ens) {Ensembles\\hard / soft vote};
  \node[grp, fit=(hc)(cnn)(ens)(hc |- G1.north)(ens |- G1.south), inner sep=0pt] (G3) {};
  \node[gl] at (G3.north) {Models};

  % 4. Evaluation
  \node[box=17mm, right=7mm of cnn, minimum height=23.4mm] (ev)
    {on the test set:\\[2pt]accuracy\\macro-F1 (CI)\\McNemar test\\per-class recall\\$Q$-statistic\\Grad-CAM};
  \node[grp, fit=(ev)(ev |- G3.north)(ev |- G3.south), inner sep=0pt] (G4) {};
  \node[gl] at (G4.north) {Test evaluation};

  \draw[arr] (pool.east) -- (G2.west |- pool);
  \draw[arr] (lt) -- (rb);
  \draw[arr] (G2.east |- lt) -- (G3.west |- cnn);
  \draw[arr] (G3.east |- cnn) -- (G4.west |- cnn);
  \draw[arr, dashed] (val.east) -- ++(2mm,0) |- ([yshift=-3mm]G2.south) -| node[pos=0.25, below=0.2mm, font=\tiny, text=wfink] {model selection: hyperparameters, early stopping} (G3.south);
  \draw[arr, dashed] (test.east) -- ++(1mm,0) |- ([yshift=-7.5mm]G2.south) -| node[pos=0.25, below=0.2mm, font=\tiny, text=wfink] {scored once per model} (G4.south);
\end{tikzpicture}
}
  \caption{Experimental workflow. A fixed balanced test set and a fixed balanced validation set are held out once; three training sets are derived from the remaining images. Every hyperparameter and early-stopping decision uses the validation set only (dashed), and the test set is scored once per model.}
  \label{fig:workflow}
\end{figure}

\subsection{Dataset and splits}
We use the SkyView aerial landscape dataset~\cite{skyview}: 12{,}000 RGB images of $256\times256$ pixels, 800 per class, in 15 classes (Agriculture, Airport, Beach, City, Desert, Forest, Grassland, Highway, Lake, Mountain, Parking, Port, Railway, Residential, River).
We hold out a stratified test set of 160 images per class (2{,}400 images).
From the remaining 640 images per class we draw a validation set of 96 images per class (1{,}440 images, seed~0), leaving a training pool of 544 images per class.
We verified that no file appears in more than one split and that the test set contains no byte-identical copy of a training image.

\subsection{Training distributions}
\label{sec:settings}
We derive three training sets from the pool (\cref{tab:settings}).

\hi{Balanced:} all 8{,}160 pool images.

\hi{Long-tailed:} class $i$ in alphabetical order keeps a fraction $r_i$ of its 544 pool images, with $r = (800, 700, 650, \ldots, 100, 50)/800$, i.e.\ from 544 images for Agriculture down to 34 for River.
The imbalance ratio is 16 and the set has 4{,}114 images.

\hi{Rebalanced:} the long-tailed set plus offline augmented copies of each minority class, generated only from that class's own long-tailed images, until every class has 544 images (4{,}046 synthetic images).
Each copy applies a random rotation in $[-30^\circ, 30^\circ]$, Gaussian blur with radius in $[0.1, 2]$, and brightness, contrast and colour-saturation factors in $[0.8, 1.2]$.

Validation and test sets stay balanced and identical in all settings.
A balanced validation set makes model selection target macro-averaged performance and keeps selection noise identical across settings, so the long-tailed and rebalanced columns differ from the balanced column only in the training data.

\begin{table}[t]
  \centering
  \caption{Training sets. Validation (96/class) and test (160/class) sets are shared and balanced.}
  \label{tab:settings}
  \begin{tabular}{@{}lrrr@{}}
    \toprule
    Setting & Images & Per class & Synthetic \\
    \midrule
    Balanced    & 8{,}160 & 544      & 0 \\
    Long-tailed & 4{,}114 & 544--34  & 0 \\
    Rebalanced  & 8{,}160 & 544      & 4{,}046 \\
    \bottomrule
  \end{tabular}
\end{table}

\subsection{Metrics and statistical testing}
We report accuracy and macro-averaged F1 on the test set; on a balanced test set accuracy equals macro-recall.
For imbalance we also report per-class recall and its mean over the five head, five middle and five tail classes of the long-tailed ordering.
Uncertainty is quantified with 95\% percentile bootstrap intervals over test images (2{,}000 resamples).
Paired comparisons between two classifiers on the same test images use the exact McNemar test~\cite{dietterich1998tests}, and we call a difference significant at $p<0.05$.
All results come from a single training run per configuration (seed~0); the bootstrap intervals capture test-sampling variability but not training variability.
